\documentclass[10pt,a4paper]{amsart}
\usepackage[margin=19mm]{geometry}
\usepackage{amsmath,amssymb,mathtools}
\usepackage[scr=rsfs,cal=euler]{mathalfa}
\usepackage{xfrac}
\usepackage[unicode,hidelinks,bookmarks=false]{hyperref}
\newcommand{\ie}{i.e.\ }
\newcommand{\cf}{cf.\ }
\newcommand{\resp}{respectively\ }
\newcommand{\Subsection}{\subsection}
\providecommand{\nobreakdash}{\nobreak\hbox{-}}
\setlength{\emergencystretch}{2em}
\multlinegap0pt
\usepackage{amssymb}
\usepackage{mathtools}
\usepackage[shortlabels]{enumitem}
\usepackage{xcolor}
\newtheorem{lemma}{Lemma}
\newcommand{\Id}{\operatorname{Id}}
\newcommand{\R}{\mathbb R}
\newcommand{\supp}{\operatorname{supp}}
\title{The maximum number of Pareto eigenvalues of a real matrix of order four is 23}
\author{Samir Adly}
\date{Source: arXiv:2607.29580v1 (CC BY 4.0)}
\begin{document}
\maketitle

\noindent\emph{Source and licence.} The maximum number of Pareto eigenvalues of a real matrix of order four is 23. Version arXiv:2607.29580v1 (CC BY 4.0), distributed by arXiv under the Creative Commons Attribution 4.0 International licence, http://creativecommons.org/licenses/by/4.0/, as recorded in the arXiv licence metadata for this exact version and shown on the abstract page https://arxiv.org/abs/2607.29580v1. Full licence notice: \texttt{licenses/arxiv-2607.29580-CC-BY-4.0.txt}.\par\smallskip

\noindent\textit{Editorial scope.} Counted unit: the article's lemma lem:simultaneous-regularization (source0.tex lines 2187--2202). The article's complete original proof of it is delivered with it (source0.tex lines 2460--2673, one \texttt{proof} environment of 214 lines, which the article places away from the statement). No mathematics is written, repaired or re-derived: the statement and the proof are the article's own lines, in the article's own notation, and no retained line is substituted or re-flowed.  No further source-local context is needed: every cross-reference of every retained line resolves inside the two delivered spans.  Nothing was omitted from any retained span, so the package carries no omission record.  No retained line contains a citation, so the wrapper carries no bibliography and no bibliography entry.  No retained line contains a figure environment, an image-inclusion command or drawing code, so no image is delivered and the package declares no asset dependency.  Editorial formatting only, in addition: an AMS \texttt{amsart} preamble that restates the article's own declarations and macros used by the retained lines, namely the environments \texttt{lemma} on separate counters, together with the standard packages those lines need.  The retained environments print in this document as Lemma 1 in order of appearance, whereas the article prints the same items with the numbering of its own full document.  The complete native source of the article as distributed in the arXiv e-print for this exact version is bundled unchanged as \texttt{sources/206501/source0.tex}.
\begin{lemma}
\label{lem:simultaneous-regularization}
Let $\lambda_1,\ldots,\lambda_k$ be distinct Pareto eigenvalues of
$A\in\R^{4\times4}$, where $1\leq k\leq3$. For each $\ell$, let
$J_\ell$ produce $\lambda_\ell$ through $u_\ell>0$. Then every
neighborhood of $A$ contains a matrix $B$ and pairwise distinct real
numbers $\beta_1,\ldots,\beta_k$ such that
$$
B_{J_\ell}u_\ell=\beta_\ell u_\ell,
\qquad
B_{J_\ell^{\mathrm c},J_\ell}u_\ell>0,
$$
and $\beta_\ell$ is an algebraically simple eigenvalue of
$B_{J_\ell}$ for every $\ell$. Thus all selected productions become
regular through the same vectors.
\end{lemma}

\begin{proof}[Proof of Lemma~\ref{lem:simultaneous-regularization}]
Extend each $u_\ell$ by zero to a vector $x_\ell\in\R^4$. No two of
these vectors are proportional, since their associated Pareto values
are distinct.

Suppose first that $x_1,\ldots,x_k$ are linearly independent. For
$J_\ell=\supp x_\ell$, choose $v_\ell$ equal to zero on $J_\ell$ and
strictly positive on $J_\ell^{\mathrm c}$. There is a linear map $H$
such that
$$
Hx_\ell=v_\ell
\qquad
(1\leq\ell\leq k).
$$
For $t>0$ small, all selected exterior slacks are strict in
$$
A_0:=A+tH,
$$
while the principal eigenvalue equations are unchanged.

Choose vectors $\ell_r\in\R^4$ such that
$$
\ell_r^\top x_s=\delta_{rs},
$$
where $\delta_{rs}$ is the Kronecker symbol, and put
$P_r=x_r\ell_r^\top$. On
$$
\R^{J_r}
=
\R x_r\oplus\bigl(\ker\ell_r^\top\cap\R^{J_r}\bigr),
$$
the matrices $(A_0)_{J_r}$ and $(P_r)_{J_r}$ have the forms
$$
\begin{pmatrix}
\lambda_r&*\\
0&C_r
\end{pmatrix},
\qquad
\begin{pmatrix}
1&0\\
0&0
\end{pmatrix}.
$$
For every sufficiently small coefficient outside a finite exceptional
set, adding this multiple of $P_r$ separates the eigenvalue associated
with $x_r$ from the quotient spectrum. Choose the coefficients
$\varepsilon_r$ successively. At each step, avoid these exceptional
values and take $\varepsilon_r$ small enough to preserve the simplicity
obtained at the preceding steps. Since
$$
P_rx_s=\delta_{rs}x_r,
$$
for the matrix
$$
B:=A_0+\sum_{r=1}^k\varepsilon_rP_r,
$$
put
$$
\beta_r:=\lambda_r+\varepsilon_r
\qquad
(1\leq r\leq k).
$$
Then $u_r$ remains an eigenvector of $B_{J_r}$, the principal
eigenvalue $\beta_r$ is simple, and the strict exterior slack is
unchanged. The coefficients may also be chosen so that the values
$\beta_1,\ldots,\beta_k$ remain pairwise distinct.

It remains to consider $k=3$ when the vectors are linearly dependent.
Since no two vectors are proportional, every coefficient in a
nontrivial dependence relation is nonzero. Because the vectors are
nonnegative, we may relabel them so that
$$
z=ax+by,
\qquad
a,b>0.
$$
Write $X=\supp x$, $Y=\supp y$, and
$Z=X\cup Y=\supp z$. The production equations give
$$
b(Ay)_i=a(\lambda_z-\lambda_x)x_i
\quad\text{for }i\in X\setminus Y,
$$
$$
a(Ax)_i=b(\lambda_z-\lambda_y)y_i
\quad\text{for }i\in Y\setminus X,
$$
and
$$
a(\lambda_x-\lambda_z)x_i
+
b(\lambda_y-\lambda_z)y_i=0
\quad\text{for }i\in X\cap Y.
$$
If $X\cap Y$, $X\setminus Y$, and $Y\setminus X$ are all nonempty,
the first two relations give
$$
\lambda_z\geq\lambda_x,
\qquad
\lambda_z\geq\lambda_y.
$$
Both inequalities are strict because the three values are distinct,
and this contradicts the third relation. Equal supports are also
impossible, since eigenvectors associated with three distinct
eigenvalues of the same principal matrix are linearly independent.
Thus the supports are disjoint or properly nested. The same relations
give, respectively,
$$
X\cap Y=\varnothing,
\qquad
\lambda_z>\max\{\lambda_x,\lambda_y\},
$$
or, after exchanging $x$ and $y$,
$$
Y\subsetneq X=Z,
\qquad
\lambda_x<\lambda_z<\lambda_y,
$$
and the restriction of $x$ to $Y$ belongs to $\R y$.

In the disjoint case, the first two relations make the slacks of $y$
on $X$ and of $x$ on $Y$ strict. In the nested case, the first
relation makes the slack of $y$ on $X\setminus Y$ strict. Thus, in
both configurations, only slacks outside $Z$ can vanish. Choose
$v_x,v_y$ equal to zero on $Z$ and positive on $Z^{\mathrm c}$, and
define
$$
Hx=v_x,
\qquad
Hy=v_y.
$$
Then $Hz=av_x+bv_y$. For $t>0$ small, the matrix
$$
A_0:=A+tH
$$
makes all three exterior slacks strict and preserves the selected
principal eigenvalue equations.

Set $W=\operatorname{span}\{x,y\}$. We construct a projection $P$ onto
$W$ which preserves the selected coordinate subspaces. If
$X\cap Y=\varnothing$, choose vectors
$\ell_x,\ell_y\in\R^4$, supported on $X,Y$, such that
$$
\ell_x^\top x=\ell_y^\top y=1,
$$
and put
$$
P=x\ell_x^\top+y\ell_y^\top.
$$
The relations above, which are unchanged for $A_0$ on $Z$, give
$$
(A_0)_Zx
=
\lambda_xx+\frac ba(\lambda_z-\lambda_y)y,
\qquad
(A_0)_Zy
=
\frac ab(\lambda_z-\lambda_x)x+\lambda_yy.
$$
Thus $W$ is invariant. Its two eigenvalues are $\lambda_z$ and
$\lambda_x+\lambda_y-\lambda_z$, which are distinct.

If $Y\subsetneq X$, choose a complement $L$ of $\R y$ in $\R^Y$.
Since $W\cap\R^Y=\R y$, extend $L$ to a complement $N_X$ of $W$ in
$\R^X$, and let $P$ be the projection onto $W$ along
$$
N_X\oplus\R^{X^{\mathrm c}}.
$$

For every selected support $J$, put $W_J=W\cap\R^J$. The active spaces
are
$$
W_X=\R x,\quad W_Y=\R y,\quad W_Z=W
$$
in the disjoint case, and
$$
W_Y=\R y,\quad W_X=W_Z=W
$$
in the nested case. Moreover,
$$
\R^J=W_J\oplus(\ker P\cap\R^J).
$$
The space $W_J$ is invariant under $(A_0)_J$, and every selected
eigenvalue is simple in the restriction to $W_J$. Relative to this
decomposition,
$$
(A_0)_J=
\begin{pmatrix}
R_J&*\\
0&Q_J
\end{pmatrix},
\qquad
P_J=
\begin{pmatrix}
\Id_{W_J}&0\\
0&0
\end{pmatrix}.
$$
A small scalar $s$, outside a finite set, separates every selected
eigenvalue of $R_J+s\Id_{W_J}$ from the spectrum of $Q_J$. Therefore
all selected principal eigenvalues of
$$
B:=A_0+sP
$$
are simple. Since $P$ fixes $x,y,z$, the three produced values are
$$
\beta_x=\lambda_x+s,
\qquad
\beta_y=\lambda_y+s,
\qquad
\beta_z=\lambda_z+s.
$$
They remain distinct and keep their strict exterior slacks. All
perturbations can be chosen in the prescribed neighborhood of $A$.
\end{proof}

\end{document}
